\documentclass[11pt]{article}
\usepackage[margin=1in]{geometry}
\usepackage{amsmath,amssymb,amsthm}
\usepackage{graphicx}
\usepackage[hidelinks]{hyperref}

\newtheorem{theorem}{Theorem}
\newtheorem{corollary}{Corollary}
\newtheorem{lemma}{Lemma}

\newcommand{\R}{\mathbb R}
\newcommand{\C}{\mathbb C}
\newcommand{\tr}{\operatorname{tr}}
\newcommand{\supp}{\operatorname{supp}}

\title{A Compact Block-Diagonal Subcase of Tracial Joint Spectral Measures}
\author{FA/Banach automatic math run}
\date{2026-07-03}

\begin{document}
\maketitle

\section*{Source And Target}

The source is Otte Heinavaara, \emph{Tracial joint spectral measures},
arXiv:2310.03227.  For Hermitian matrices $A,B$, the source constructs a
positive measure $\mu_{A,B}$ on $\R^2$ satisfying
\begin{equation}\label{eq:source}
  \tr H(f)(xA+yB)=\int_{\R^2} f(ax+by)\,d\mu_{A,B}(a,b),
\end{equation}
where
\[
  H(f)(s)=\int_0^1 \frac{1-t}{t} f(ts)\,dt .
\]
Section 1.8 of the source lists as future work the ``existence and uniqueness
of the measures $\mu_{A,B}$ for compact self-adjoint operators $A$ and $B$.''

\begin{center}
\includegraphics[width=0.95\linewidth]{figures/open_problem_crop.png}
\end{center}

\section*{Result}

We prove a compact-operator subcase: common finite-dimensional reducing blocks
with square-summable Hilbert-Schmidt size.  This includes the important
commuting compact Hilbert-Schmidt case, by simultaneous diagonalization.

\section*{Idea Of The Proof}

The finite-dimensional measure is already available on each reducing block.
Because the blocks are square-summable in Hilbert-Schmidt norm, the block
measures have summable second moments.  This gives local finiteness away from
the origin, so their countable sum is a Radon measure on the punctured plane.
The source identity then sums block by block for tests supported away from
zero.  Uniqueness follows from the same moment-determinacy mechanism used in
the source: after multiplying by $a^2+b^2$, one obtains a finite compactly
supported measure whose polynomial moments are fixed.

\begin{theorem}[Block compact tracial joint spectral measures]\label{thm:block}
Let $\mathcal H=\bigoplus_{j=1}^\infty \mathcal H_j$ be an orthogonal direct
sum of finite-dimensional Hilbert spaces.  Let
\[
  A=\bigoplus_{j=1}^\infty A_j,\qquad B=\bigoplus_{j=1}^\infty B_j
\]
be self-adjoint block-diagonal compact operators, where
$A_j,B_j\in B(\mathcal H_j)$ are Hermitian matrices, and assume
\begin{equation}\label{eq:hs}
  \sum_{j=1}^\infty \bigl(\tr A_j^2+\tr B_j^2\bigr)<\infty .
\end{equation}
For each $j$, let $\mu_j=\mu_{A_j,B_j}$ be the finite-dimensional tracial
joint spectral measure of the source.  Then
\[
  \mu_{A,B}:=\sum_{j=1}^\infty \mu_j
\]
defines a positive Radon measure on $\R^2\setminus\{0\}$, supported in
$\{(a,b): |a|+|b|\le \|A\|+\|B\|\}$, with finite second moment
\begin{equation}\label{eq:second}
  \int_{\R^2\setminus\{0\}} (a^2+b^2)\,d\mu_{A,B}(a,b)
  =\frac{1}{6}\sum_{j=1}^\infty \bigl(\tr A_j^2+\tr B_j^2\bigr).
\end{equation}
Moreover, for every bounded Borel $f:\R\to\R$ whose support is compact and
does not contain $0$, and every $x,y\in\R$,
\begin{equation}\label{eq:identity}
  \tr H(f)(xA+yB)=
  \int_{\R^2\setminus\{0\}} f(ax+by)\,d\mu_{A,B}(a,b).
\end{equation}
For every $p\ge 2$,
\begin{equation}\label{eq:p}
  \tr |xA+yB|^p
  =p(p+1)\int_{\R^2\setminus\{0\}} |ax+by|^p\,d\mu_{A,B}(a,b).
\end{equation}
Finally, $\mu_{A,B}$ is unique among positive Radon measures on
$\R^2\setminus\{0\}$ with bounded support and finite second moment satisfying
\eqref{eq:identity} and the polynomial identities obtained from
\eqref{eq:p} for integer $p\ge 2$.
\end{theorem}

\begin{proof}
For a finite-dimensional Hermitian pair $(C,D)$, apply the source identity
with $f(t)=t^2$.  Since $H(f)(t)=t^2/6$, polarization gives
\[
  \int_{\R^2} (ax+by)^2\,d\mu_{C,D}(a,b)
  =\frac{1}{6}\tr(xC+yD)^2 .
\]
Taking $(x,y)=(1,0),(0,1)$ and adding yields
\begin{equation}\label{eq:block-second}
  \int_{\R^2}(a^2+b^2)\,d\mu_{C,D}(a,b)
  =\frac{1}{6}\bigl(\tr C^2+\tr D^2\bigr).
\end{equation}
The source uniqueness proof also gives the support bound
$\supp\mu_{C,D}\subseteq\{(a,b): |a|+|b|\le \|C\|+\|D\|\}$.

Let $K\subset \R^2\setminus\{0\}$ be compact and put
$\delta=\operatorname{dist}(K,0)>0$.  By \eqref{eq:block-second},
\[
  \sum_j \mu_j(K)
  \le \delta^{-2}\sum_j \int_K(a^2+b^2)\,d\mu_j
  \le \frac{1}{6\delta^2}\sum_j(\tr A_j^2+\tr B_j^2)<\infty .
\]
Thus $\sum_j\mu_j$ is locally finite on the punctured plane.  The same
calculation proves the second-moment formula \eqref{eq:second}.  The support
bound follows from $\|A_j\|\le \|A\|$ and $\|B_j\|\le \|B\|$.

Now let $f$ be bounded and supported in a compact subset of
$\R\setminus\{0\}$.  For fixed $(x,y)\ne(0,0)$, the function
$(a,b)\mapsto f(ax+by)$ is supported in a closed subset of the bounded support
of $\mu_{A,B}$ that stays a positive distance from the origin.  Hence the
right side of \eqref{eq:identity} is the absolutely convergent sum of the block
integrals.  On each block, the source identity gives
\[
  \tr H(f)(xA_j+yB_j)
  =\int_{\R^2} f(ax+by)\,d\mu_j(a,b).
\]
Summing in $j$ proves \eqref{eq:identity}; the absolute convergence also shows
that $H(f)(xA+yB)$ is trace class.  The case $(x,y)=(0,0)$ is immediate.

For $p\ge2$, finite-dimensional source identity with $f(t)=|t|^p$ gives the
block formula.  Since
\[
  \tr|xA_j+yB_j|^p\le
  \|xA+yB\|^{p-2}\tr(xA_j+yB_j)^2
\]
and the right side is summable by \eqref{eq:hs}, the block traces sum
absolutely.  Summing the block identities gives \eqref{eq:p}.

It remains to justify uniqueness.  Suppose $\nu$ is another positive Radon
measure on $\R^2\setminus\{0\}$ with bounded support and finite second moment,
satisfying the same polynomial identities.  Define finite compactly supported
measures
\[
  d\widetilde\mu=(a^2+b^2)\,d\mu_{A,B},\qquad
  d\widetilde\nu=(a^2+b^2)\,d\nu .
\]
For every homogeneous polynomial $q$ of degree at least $2$, polarization of
the identities for $(ax+by)^m$ gives
\[
  \int q(a,b)\,d\mu_{A,B}=\int q(a,b)\,d\nu .
\]
Consequently, for every polynomial $r(a,b)$,
\[
  \int r(a,b)\,d\widetilde\mu
  =\int r(a,b)(a^2+b^2)\,d\mu_{A,B}
  =\int r(a,b)(a^2+b^2)\,d\nu
  =\int r(a,b)\,d\widetilde\nu .
\]
Compactly supported finite Borel measures on $\R^2$ are determined by their
polynomial moments, so $\widetilde\mu=\widetilde\nu$.  Dividing by $a^2+b^2$
on the punctured plane gives $\mu_{A,B}=\nu$.
\end{proof}

\begin{corollary}[Commuting compact Hilbert-Schmidt pairs]\label{cor:commuting}
Let $A$ and $B$ be commuting compact self-adjoint Hilbert-Schmidt operators on
a separable Hilbert space.  Then $A$ and $B$ admit a unique tracial joint
spectral measure in the sense of Theorem~\ref{thm:block}.
\end{corollary}

\begin{proof}
Commuting compact self-adjoint operators are simultaneously diagonalizable:
there is an orthonormal basis in which $A=\operatorname{diag}(a_j)$ and
$B=\operatorname{diag}(b_j)$.  The Hilbert-Schmidt assumption says
$\sum_j(a_j^2+b_j^2)<\infty$, so Theorem~\ref{thm:block} applies with
one-dimensional blocks.  In this case the measure is explicit:
\[
  \mu_{A,B}=\sum_j (t\mapsto t(a_j,b_j))_\#
  \left(\frac{1-t}{t}\,dt\bigg|_{(0,1]}\right).
\]
\end{proof}

\section*{Verification Notes}

The script \texttt{code/verify\_commuting\_model.py} checks the normalization
in Corollary~\ref{cor:commuting}.  For a diagonal Hilbert-Schmidt model it
verifies
\[
  \int_0^1 \frac{1-t}{t}|t(xa_j+yb_j)|^p\,dt
  =\frac{|xa_j+yb_j|^p}{p(p+1)}
\]
and prints the finite second-moment matrix.

\section*{Limitations And Novelty Check}

This is a partial result.  It does not prove existence for arbitrary compact
self-adjoint pairs.  A natural compression proof for all Hilbert-Schmidt pairs
would require a continuity/tightness theorem for finite-dimensional tracial
joint spectral measures that prevents second-moment loss at the origin; that
step is not proved here.

Cheap run indexes were searched for arXiv:2310.03227 and for tracial joint
spectral measure keywords.  The source bibliography lists the follow-up
``Properties of Tracial Joint Spectral Measures'' as in preparation, and a
bounded web search for that title and for compact self-adjoint tracial joint
spectral measures did not find an available arXiv answer.

\section*{References}

\begin{thebibliography}{9}
\bibitem{HeinavaaraTJSM}
O. Heinavaara, \emph{Tracial joint spectral measures}, arXiv:2310.03227,
2023.
\end{thebibliography}

\end{document}
